\section{Open questions}\label{sec:open}

\paragraph{Earth Mover's Distance.}
The Earth Mover's Distance (EMD) between two planar point sets sits in the same range-counting framework of
Driemel et al.~\cite{DMOSW25} as MST weight, and there too a polynomial gap separates the known upper
bound from the lower bound. Sublinear and sketching algorithms for EMD have a long
history~\cite{AndoniDoBaIndykWoodruff09,DoBaEtal11,Charikar02}, as do near-linear geometric matching
algorithms~\cite{AgarwalChangRaghvendraXiao22,AgarwalChangXiao19}, but none operates under a pure
range-counting oracle.

We believe the empty-cell spatial leader estimator of \cref{sec:support-estimator} is the right primitive
here, and sketch why. The estimator is not specific to single-linkage clustering: at its core it estimates
the number of connected components of an \emph{implicitly defined} geometric graph from a sample of cells
drawn from a spatial universe that includes empty cells, with variance governed by the component count and
not the occupied-cell count. EMD has a dyadic structure parallel to the cluster-count integral. Writing
the two point sets as red and blue with equal mass, the cost of an optimal transport admits a quadtree
approximation: at each dyadic level one pays, per cell, the absolute red-minus-blue mass imbalance times
the cell diameter, and summing over levels approximates EMD within an $O(\log\Delta)$ factor
(the source of the multiplicative $\log\Delta$ loss already present in~\cite{DMOSW25}). The analogue of
the cluster-count curve is thus the imbalance profile across scales, and the analogue of estimating
$c_S(t)$ at a scale is estimating the number of cells with nonzero imbalance, or more precisely the total
imbalance, at that scale.

Two ingredients would need EMD analogues. First, a \emph{spatial estimator} for the per-scale imbalance:
the leader estimator counts components, while EMD needs a (signed) mass functional, so one would draw a
candidate cell uniformly and return its red-minus-blue count scaled by the universe size, an unbiased
estimator whose variance is again controlled by a spatial, not mass-proportional, sample; the heavy-tail
difficulty (a few cells carrying most of the imbalance) is exactly the one the empty-cell view was built
for. Second, a \emph{packing inequality} relating the number of imbalanced cells at a scale to the EMD
itself, playing the role of $\Wq\ge b\delta/16$ and collapsing the geometric sum of per-scale costs.
Establishing these two and balancing the resulting $\sqrt K+n/K$ trade-off would, we expect, bring planar
EMD to the same $\Otil(n^{1/3})$ rate, matching the lower bound of~\cite{DMOSW25} up to the unavoidable
$\log\Delta$ transport-approximation factor. Whether the spatial-sampling idea extends to the full
matching cost, and not only to its dyadic approximation, is the deeper question.

\paragraph{Higher dimensions.}
For $P\subseteq[\Delta]^d$ with $d\ge3$ the WSPD spanner and the cluster-count integral both go through,
but the packing exponent and the balance point $K=\Theta(n^{2/3})$ are tuned to the plane. The
dimension-dependence of the optimal range-counting exponent, and whether it degrades gracefully with $d$,
is open.

\paragraph{Polylogarithmic factors and determinism.}
\Cref{cor:tight} leaves a polylogarithmic and $\eps$-dependent gap. Pinning down the exact dependence on
$\eps$ and $\log\Delta$, and asking whether a deterministic range-counting algorithm can match the
$n^{1/3}$ rate or whether randomization is necessary at this oracle, are concrete questions about the
constant in the exponent's neighborhood.
